% ================================================================
% PREAMBLE AND TYPOGRAPHICAL CORRECTIONS
% ================================================================

% Add to the theorem declarations in the preamble:
\newtheorem{corollary}[theorem]{Corollary}

% Correct the Section 2 title to:
% \section{Graph-Laplacian formulation of adaptive resistor-network learning}

% ================================================================
% REPLACE THE NULL-SPACE NON-INTERFERENCE THEOREM IN SECTION 4
% ================================================================

\begin{theorem}[Sharp null-space non-interference]
Assume that the old task is fitted exactly at \(g_A\), so that
\[
    r_A(g_A)=0.
\]
If a conductance perturbation \(\Delta g\) satisfies
\[
    J_A(g_A)\Delta g=0,
\]
then
\[
    F_A(\Delta g)
    =
    \frac{1}{8}
    \left\|
    D^2y_A(g_A)[\Delta g,\Delta g]
    \right\|^2
    +
    O(\|\Delta g\|^5).
\]
In particular,
\[
    \boxed{
    F_A(\Delta g)=O(\|\Delta g\|^4).
    }
\]
Thus, an update in \(\ker J_A(g_A)\) removes not only the first-order
output distortion but also the quadratic-order contribution to the old-task
loss.
\end{theorem}

\begin{proof}
Since \(r_A(g_A)=0\), Taylor expansion of the residual gives
\[
    r_A(g_A+\Delta g)
    =
    J_A(g_A)\Delta g
    +
    \frac{1}{2}D^2y_A(g_A)[\Delta g,\Delta g]
    +
    O(\|\Delta g\|^3).
\]
If \(J_A(g_A)\Delta g=0\), then
\[
    r_A(g_A+\Delta g)
    =
    \frac{1}{2}D^2y_A(g_A)[\Delta g,\Delta g]
    +
    O(\|\Delta g\|^3).
\]
Because
\[
    F_A(\Delta g)
    =
    \frac{1}{2}
    \|r_A(g_A+\Delta g)\|^2,
\]
we obtain
\[
    F_A(\Delta g)
    =
    \frac{1}{8}
    \left\|
    D^2y_A(g_A)[\Delta g,\Delta g]
    \right\|^2
    +
    O(\|\Delta g\|^5).
\]
The leading term is fourth order in \(\Delta g\), which proves the result.
\end{proof}

% ================================================================
% REPLACE THE CORRESPONDING SENTENCE IN THE NUMERICAL SECTION
% ================================================================

% Replace:
%   F_A(\Delta g_\perp)=O(\varepsilon^3)
% with:
\[
    F_A(\Delta g_\perp)=O(\varepsilon^4),
\]
% and state that the null-space log--log slope should approach 4.

% ================================================================
% ADD TO THE DISCUSSION OF SENSITIVITY ANCHORING
% ================================================================

Although sensitivity anchoring removes the quadratic forgetting term in the
large-\(\lambda\) limit, the limiting safe update need not vanish. Consequently,
the full nonlinear forgetting generally approaches a quartic residual of the
form
\[
    F_A(\Delta g_{\mathrm{safe}})
    =
    \frac{1}{8}
    \left\|
    D^2y_A(g_A)
    [\Delta g_{\mathrm{safe}},\Delta g_{\mathrm{safe}}]
    \right\|^2
    +
    O(\|\Delta g_{\mathrm{safe}}\|^5),
\]
rather than converging identically to zero. This distinguishes suppression of
quadratic-order interference from exact nonlinear retention.
